%-------------------------
% Resume in Latex
% Author : Jake Gutierrez
% Based off of: https://github.com/sb2nov/resume
% License : MIT
%------------------------

\documentclass[a4paper,11pt]{article}
\usepackage[empty]{fullpage}
\usepackage{titlesec}
\usepackage[usenames,dvipsnames]{color}
\usepackage{enumitem}
\usepackage{fancyhdr}
\usepackage[portuguese]{babel}
\usepackage{tabularx}
\input{glyphtounicode}
\usepackage[document]{ragged2e}
\usepackage{fontawesome}
\usepackage[hidelinks]{hyperref}
\hypersetup{pdfnewwindow=true}

\pagestyle{fancy}
\setlength{\footskip}{4.5pt}
\fancyhf{} % clear all header and footer fields
\fancyfoot{}
\renewcommand{\headrulewidth}{0pt}
\renewcommand{\footrulewidth}{0pt}

% Adjust margins
\addtolength{\oddsidemargin}{-0.5in}
\addtolength{\evensidemargin}{-0.5in}
\addtolength{\textwidth}{1in}
\addtolength{\topmargin}{-.5in}
\addtolength{\textheight}{1.0in}

\urlstyle{same}

\raggedbottom
\raggedright
\setlength{\tabcolsep}{0in}

% Sections formatting
\titleformat{\section}{
  \vspace{-4pt}\scshape\raggedright\large
}{}{0em}{}[\color{black}\titlerule \vspace{-5pt}]

% Ensure that generate pdf is machine readable/ATS parsable
\pdfgentounicode=1

%-------------------------
% Custom commands
\newcommand{\resumeItem}[1]{
  \item\small{
    {#1 \vspace{-2pt}}
  }
}

\newcommand{\resumeSubheading}[4]{
  \vspace{-2pt}\item
    \begin{tabular*}{0.97\textwidth}[t]{l@{\extracolsep{\fill}}r}
      \textbf{#1} & #2 \\
      \textit{\small#3} & \textit{\small #4} \\
    \end{tabular*}\vspace{-7pt}
}

\newcommand{\resumeSubSubheading}[2]{
    \item
    \begin{tabular*}{0.97\textwidth}{l@{\extracolsep{\fill}}r}
      \textit{\small#1} & \textit{\small #2} \\
    \end{tabular*}\vspace{-7pt}
}

\newcommand{\resumeProjectHeading}[2]{
    \item
    \begin{tabular*}{0.97\textwidth}{l@{\extracolsep{\fill}}r}
      \small#1 & #2 \\
    \end{tabular*}\vspace{-7pt}
}

\newcommand{\resumeSubItem}[1]{\resumeItem{#1}\vspace{-4pt}}

\renewcommand\labelitemii{$\vcenter{\hbox{\tiny$\bullet$}}$}

\newcommand{\resumeSubHeadingListStart}{\begin{itemize}[leftmargin=0.15in, label={}]}
\newcommand{\resumeSubHeadingListEnd}{\end{itemize}}
\newcommand{\resumeItemListStart}{\begin{itemize}}
\newcommand{\resumeItemListEnd}{\end{itemize}\vspace{-5pt}}

%-------------------------------------------
%%%%%%  RESUME STARTS HERE  %%%%%%%%%%%%%%%%%%%%%%%%%%%%


\begin{document}

\begin{center}
    {\huge\textsc{Pedro Primor}} \\ \vspace{1pt}
    \ \href{mailto:hello@primor.me}{\faEnvelope \hspace{1pt} hello@primor.me}
    % $|$ \faLink \ \href{https://primor.me/}{\underline{primor.me}}
    $|$ \href{https://www.linkedin.com/in/primor/}{\faLinkedin/primor}
    $|$ \href{https://github.com/pprimor}{\faGithub/pprimor}
\end{center}

\section{Sobre mim}
    \justifying

Sou engenheiro de software na \href{https://pageproof.com/}{PageProof}, uma pequena equipa remota que faz revisão e aprovação de conteúdos online para equipas de marketing e criativas. A minha parte é a plataforma de integrações: os plugins que permitem rever provas dentro do Canva, Microsoft Office, Adobe Creative Cloud, Final Cut Pro e Jira, o servidor OAuth 2.0 / OpenID Connect que os autentica e a aplicação Electron para computador. Quase tudo é TypeScript, React e Node.js; Swift e C++ entram quando a aplicação anfitriã não deixa outra opção. \textbf{Idiomas:} português (nativo) e inglês (fluente).

%-----------SKILLS-----------
\section{Competências}
  \resumeSubHeadingListStart
    \item\small{
      \textbf{Linguagens:} TypeScript, JavaScript, Node.js, Swift, C++ \\
      \textbf{Frameworks e plataformas:} React, Electron, Vite, Adobe UXP, Canva Apps SDK, Office.js, Atlassian Forge \\
      \textbf{Infraestrutura:} OAuth 2.0 / OIDC, Kubernetes, Azure, Cloudflare Workers, Redis, GitHub Actions
    }\vspace{-5pt}
  \resumeSubHeadingListEnd

%-----------EXPERIENCE-----------
\section{Experiência profissional}
  \resumeSubHeadingListStart

    \resumeSubheading
      {Engenheiro de Software}{mar. 2023 -- atual}
      {PageProof.com Limited}{Remoto}
        \resumeItemListStart
            \resumeItem{Quem revê trabalha dentro das ferramentas de desenho e edição, por isso as \href{https://pageproof.com/integrations}{integrações} levam a revisão até lá. Construí os plugins para Canva, Microsoft Office, Final Cut Pro e QuarkXPress, co-desenvolvi os plugins Adobe UXP e construí o painel do Premiere Pro, que ordena os comentários por \emph{timecode} para o editor saltar directamente para cada um.}
            \resumeItem{A integração com o Jira liga uma prova à sua \emph{issue} e mantém as duas a par: comentários de estado, ligações remotas e o ficheiro aprovado sincronizam automaticamente, num serviço Node.js com \emph{locks} em Redis para manter os eventos em ordem. A aplicação Atlassian Forge está em testes finais.}
            \resumeItem{Todos os plugins e a aplicação para computador precisam de saber quem é o utilizador, por isso concebi um único servidor de autorização OAuth 2.0 (\href{https://www.rfc-editor.org/rfc/rfc6749.html}{RFC 6749}) com OpenID Connect, em Node.js, que todos partilham. Está agora a migrar para Kubernetes.}
            \resumeItem{Sou o principal engenheiro da aplicação Electron para computador. O seu painel de comentários foi extraído para um pacote partilhado com os plugins Adobe, para que uma correcção chegue a todas as aplicações de uma vez.}
            \resumeItem{O suplemento do Office faz agora verificações assistidas por IA antes de um ficheiro se tornar prova: imagens dos diapositivos no PowerPoint, PDFs posicionados no Word e Excel. A entrega de ficheiros aprovados para Google Drive, Dropbox e OneDrive está em curso.}
        \resumeItemListEnd

    \resumeSubheading
      {Analista de Investigação Operacional}{jan. 2022 -- set. 2022}
      {CTT – Correios de Portugal, S.A.}{Lisboa, Portugal}
        \resumeItemListStart
            \resumeItem{Trabalhei na última etapa da distribuição do correio, em que o carteiro sai do centro de distribuição e visita cada código postal do seu giro. O objectivo era um modelo de rotas que os CTT pudessem aplicar a qualquer centro do país.}
            \resumeItem{Modelei-o como um problema de rotas de veículos com as restrições dos CTT e agrupei códigos postais em segmentos para as rotas se manterem estáveis de um dia para o outro, com \emph{dashboards} em Tableau para o planeamento. Este trabalho deu origem à minha tese de mestrado.}
        \resumeItemListEnd

  \resumeSubHeadingListEnd

%-----------PROJECTS-----------
\section{Projetos}
  \resumeSubHeadingListStart

    \resumeProjectHeading
      {\textbf{\href{https://github.com/pprimor/igcp-aforro}{igcp-aforro}} $|$ \emph{biblioteca TypeScript \& CLI}}{}
      \resumeItemListStart
        \resumeItem{Um simulador de Certificados de Aforro (IGCP, Séries B a F): aritmética decimal exacta, entradas validadas com Zod e testes de referência contra as taxas que o IGCP publica.}
        \resumeItem{Publicado no npm com CLI, um \texttt{rates.json} estático, uma API HTTP só de leitura e um servidor MCP; a CI actualiza os dados Euribor e IGCP.}
      \resumeItemListEnd

  \resumeSubHeadingListEnd

%-----------EDUCATION-----------
\section{Educação}
\resumeSubHeadingListStart

  \resumeSubheading
    {Universidade NOVA de Lisboa}{set. 2020 -- dez. 2022}
    {Mestrado em Ciência da Computação}{Lisboa, Portugal}

  \resumeSubheading
    {Technische Universität Dresden}{abr. 2021 -- out. 2021}
    {Mobilidade Erasmus+ em Ciência da Computação}{Dresden, Alemanha}

  \resumeSubheading
    {Universidade NOVA de Lisboa}{set. 2017 -- set. 2020}
    {Licenciatura em Ciência da Computação}{Lisboa, Portugal}

\resumeSubHeadingListEnd

\end{document}
